\section{Related work}
\label{sec:related}
% Sources: Wang 2019 abstract ("outperforms a single-prediction baseline and dropout-based
% multiple predictions"; "helps to reduce overconfident incorrect predictions"); Moshkov 2020
% conclusions; Shanmugam 2021 abstract; Jungo & Reyes 2019 abstract; Zenk 2025 abstract;
% Fan 2020 Tables 1-2 (ClinicDB 0.899, Kvasir 0.898; unseen ColonDB 0.709, ETIS 0.628).

\paragraph{TTA in medical segmentation.} \citet{ayhan2018tta} and \citet{wang2019aleatoric} used TTA to estimate aleatoric uncertainty in medical classification and segmentation, the latter reporting that it outperformed a single prediction and dropout-based multiple predictions and reduced overconfident incorrect predictions. \citet{moshkov2020tta} reported higher average accuracy with TTA for U-Net and Mask R-CNN on microscopy images. \citet{shanmugam2021better} showed for classification that TTA can turn many correct predictions into incorrect ones even when accuracy rises. We look for the segmentation counterpart: lesions that disappear while mean Dice rises. Throughout, TTA is test-time augmentation, which changes only how a fixed model is queried, and not test-time adaptation, which updates the model on the test distribution.

\paragraph{Uncertainty and failure detection.} Deep ensembles \cite{lakshminarayanan2017ensembles} and Monte Carlo dropout \cite{gal2016dropout} are the common alternatives. \citet{jungo2019reliability} found segmentation uncertainty estimates well calibrated at the dataset level but miscalibrated at the subject level. \citet{zenk2025failure} benchmarked failure detection for 3D segmentation under realistic shifts, advocated risk--coverage analysis, and found the pairwise Dice between ensemble predictions a simple and robust baseline. We use the same risk--coverage view and separate failures with a non-empty prediction from empty predictions, which ROI-based signals cannot score.

\paragraph{Evaluation under shift.} Polyp segmentation models are commonly trained on Kvasir-SEG and CVC-ClinicDB and tested on unseen collections, where accuracy falls \cite{fan2020pranet}. PolypGen \cite{ali2023polypgen} is the purpose-built multi-centre alternative, with the centre labels and sequence structure the legacy collections lack. We used the legacy collections because published polyp models report against them, so our baselines are comparable; the cost is that we cannot group images by centre or by polyp (Section~\ref{sec:eval}). We ask a different question: whether an inference-time change that raises accuracy under shift also changes which lesions are found.

\paragraph{How misses are counted.} \citet{maierhein2024metrics} and \citet{reinke2024pitfalls} recommend instance-level detection metrics when the question is whether a lesion was found. We count images instead because the external releases carry no polyp identities, so a per-instance count is not recoverable; within that limit we report three image-level criteria side by side, which has not to our knowledge been done for TTA.
